
\documentclass[runningheads]{llncs}

\usepackage[T1]{fontenc}
\usepackage{graphicx}
\usepackage{amsmath,amssymb}
\usepackage{booktabs}
\usepackage{array}
\usepackage{multirow}
\usepackage{url}

\begin{document}

\title{A Hybrid Quantum--Classical State-Space Framework for Cyber-Defense Optimization: LQR-to-QUBO Formulation and IBM Quantum Validation}

\titlerunning{Hybrid Quantum--Classical Cyber-Defense Optimization}

\author{
Souala Elhoussine\inst{1} \and
Hadji Ibrahim\inst{1} \and
Mohamed Guerroumi\inst{1} \and
Abdelouahid Derhab\inst{2}
}

\authorrunning{S. Elhoussine et al.}

\institute{
University of Science and Technology Houari Boumediene (USTHB), Algiers, Algeria
\and
Center of Excellence in Information Assurance (CoEIA), King Saud University, Riyadh, Saudi Arabia
}

\maketitle

\begin{abstract}
Cyber-defense resource allocation is a dynamic decision problem in which a defender must protect a subset of network nodes while the attack state evolves over time. This paper presents a hybrid quantum--classical framework that formulates cyber defense as a finite-horizon state-space control problem and solves the resulting binary optimization problem using the Quantum Approximate Optimization Algorithm (QAOA) on real IBM quantum hardware. The network is modeled through an extended cyber state vector containing node compromise level, attacker knowledge, vulnerability exposure, and asset criticality. A Linear Quadratic Regulator (LQR)-style cost is expanded over a planning horizon and converted exactly into a Quadratic Unconstrained Binary Optimization (QUBO) problem. 

The evaluation is organized into two panels. The first panel validates the method on a 16-qubit instance, where Greedy, CP-SAT, Local QAOA, and IBM QAOA solve the same full QUBO. Local QAOA and IBM QAOA achieve gaps of $0.952\%$ and $1.476\%$ relative to the CP-SAT reference, respectively, while the Greedy baseline is $210.676\%$ away from CP-SAT. The second panel evaluates hardware-scale instances of 64, 128, and 152 qubits using CP-SAT and IBM QAOA. IBM QAOA remains within $4.132\%$, $7.079\%$, and $7.900\%$ of CP-SAT, respectively. These results do not establish quantum advantage; instead, they demonstrate that an LQR-derived cyber-defense QUBO can be executed on current IBM hardware while preserving solution quality close to a strong classical reference.

\keywords{Cyber Defense \and QAOA \and QUBO \and LQR \and State-Space Control \and IBM Quantum \and CP-SAT \and Hybrid Quantum-Classical Computing}
\end{abstract}

\section{Introduction}

Modern cyber defense is not a one-shot classification task. A network changes continuously as adversaries gather information, exploit vulnerable assets, and move across connected systems. A defender must therefore decide not only which nodes are currently compromised, but also which nodes should receive limited defensive resources over a future decision horizon. This is naturally a control problem: the network has a state, the defender applies an action, an attack disturbance affects the system, and a new state is produced.

Classical control theory provides a principled language for such problems. In particular, Linear Quadratic Regulator (LQR) formulations balance two competing objectives: reducing the risk represented by the system state and avoiding excessive control effort. In cyber defense, these two terms correspond to reducing compromise, attacker knowledge, vulnerability exposure, and critical-asset risk, while avoiding unnecessary or costly defense actions. Because defense decisions are often discrete---for example, patch a node or do not patch it---the resulting control problem can be formulated as a binary optimization problem.

Quantum optimization offers a different way of sampling binary solutions. The Quantum Approximate Optimization Algorithm (QAOA) prepares a parameterized quantum state over candidate bitstrings and uses alternating problem and mixer Hamiltonians to bias measurement outcomes toward low-cost solutions. QAOA is particularly suitable for Quadratic Unconstrained Binary Optimization (QUBO) problems, which are also common in classical combinatorial optimization. The central idea of this work is therefore to use LQR as a mathematically interpretable generator of cyber-defense QUBOs and to evaluate whether those QUBOs can be solved on current IBM quantum hardware.

An earlier version of this study compared uncontrolled defense, local QAOA, and hybrid IBM QAOA. While that formulation showed feasibility, the uncontrolled baseline was weak: a strong reviewer could reasonably ask whether a classical optimizer would solve the same QUBO equally well or better. The present work addresses this limitation by introducing CP-SAT as the primary classical reference and by separating the evaluation into two experimental panels. The validation panel compares Greedy, CP-SAT, Local QAOA, and IBM QAOA on the same 16-qubit full QUBO. The scalability panel compares CP-SAT and IBM QAOA on larger 64-, 128-, and 152-qubit instances. The primary metric is the gap to CP-SAT, rather than raw closed-loop cost across different problem sizes.

The contributions of this paper are:
\begin{itemize}
    \item an extended state-space cyber-defense model incorporating compromise level, attacker knowledge, vulnerability exposure, and asset criticality;
    \item an exact finite-horizon LQR-to-QUBO transformation for binary cyber-defense actions;
    \item a hybrid quantum--classical implementation using Local QAOA simulation and IBM quantum hardware;
    \item a CP-SAT benchmark methodology based on gap-to-reference evaluation;
    \item validation and hardware-scale experiments up to 152 qubits.
\end{itemize}

We emphasize that this paper does not claim quantum advantage. The supported claim is more specific and experimentally defensible: IBM QAOA produces solutions comparable to a strong classical reference on validation instances and remains within a moderate gap on larger hardware-scale instances.

\section{Related Work}

Quantum approaches to cybersecurity have been studied primarily in cryptography and machine learning. Quantum key distribution introduced the possibility of information-theoretic security for communication channels~\cite{bennett1984quantum}. More recent work has explored quantum machine learning for intrusion detection and threat intelligence~\cite{garg2022quantum}. These directions differ from the present work, which treats cyber defense as a control and optimization problem rather than a classification or communication problem.

QAOA was introduced by Farhi, Goldstone, and Gutmann as a variational quantum algorithm for approximate combinatorial optimization~\cite{farhi2014qaoa}. It alternates between a problem Hamiltonian encoding the objective and a mixer Hamiltonian that promotes exploration of the binary solution space. Variational quantum algorithms are attractive for NISQ hardware, but they also suffer from limitations such as noise sensitivity, finite sampling, and barren plateau effects~\cite{cerezo2021variational,mcclean2018barren}. These limitations motivate careful benchmarking against classical solvers.

State-space models have also appeared in security analysis, especially in epidemic spreading, positive systems, and game-theoretic moving-target defense~\cite{khanafer2016stability,zhu2013game}. Astafieva and Stepanenko emphasize controllability as a key property in optimal cybersecurity management~\cite{astafieva2025optimal}. The present paper connects these control-theoretic ideas with QUBO-based quantum optimization. The benchmarking design follows the broader view that optimization applications should be evaluated against transparent classical references and hardware-aware metrics~\cite{lubinski2024optimization}.

\section{Cyber-Defense State-Space Model}

\subsection{Extended Cyber State}

Consider a network with $n$ nodes. For each node $i$, the cyber state contains four quantities:
\begin{itemize}
    \item $s_i \in [0,1]$: compromise level of the node;
    \item $\theta_i \in [0,1]$: attacker knowledge about the node;
    \item $v_i \in [0,1]$: vulnerability exposure of the node;
    \item $c_i \in [0,1]$: asset criticality of the node.
\end{itemize}

The global state vector is
\begin{equation}
x_k =
\begin{bmatrix}
s_k\\
\theta_k\\
v_k\\
c_k
\end{bmatrix}
\in \mathbb{R}^{4n},
\end{equation}
where each block is in $\mathbb{R}^{n}$. This representation is more realistic than using compromise and attacker knowledge alone. Two nodes may have the same compromise level but different security priorities if one is a database server and the other is a low-value workstation. Similarly, a currently uncompromised node may still require attention if it is highly vulnerable or highly critical.

\subsection{Networked Dynamics}

Let $W \in \mathbb{R}^{n \times n}$ denote a row-normalized network adjacency matrix. In the implementation, $W$ is generated using a Barabási--Albert graph to introduce heterogeneous connectivity. The state evolves as
\begin{equation}
x_{k+1}=Ax_k+Bu_k+Fw_k,
\label{eq:state}
\end{equation}
where $u_k \in \{0,1\}^{n}$ is the binary defense action and $w_k \in \mathbb{R}^{n}$ is the attack disturbance.

The implemented block structure is
\begin{align}
s_{k+1} &=
(1-\delta_s)s_k +
(\beta_{\mathrm{self}}I+\beta_{\mathrm{net}}W)\theta_k +
\rho_v v_k, \\
\theta_{k+1} &=
\left((1-\gamma_\theta)I+\theta_{\mathrm{spread}}W\right)\theta_k
-\eta_\theta u_k+\alpha w_k, \\
v_{k+1} &=
(1-\lambda_v)v_k-\eta_v u_k+0.20w_k, \\
c_{k+1} &= c_k.
\end{align}

The control input reduces attacker knowledge and vulnerability exposure. The disturbance injects attacker pressure into knowledge and vulnerability. Asset criticality is modeled as slowly varying over the control horizon, so it remains constant during one optimization instance. The cost weights used in the experiments are
\begin{equation}
Q=\mathrm{diag}(4I_n,3I_n,2I_n,1.5I_n), \qquad R=0.08I_n.
\end{equation}
Thus, compromise is penalized most strongly, followed by attacker knowledge, vulnerability exposure, and asset criticality.

\section{Finite-Horizon LQR-to-QUBO Formulation}

\subsection{Finite-Horizon Control Objective}

For a horizon $H$, the defender optimizes a sequence of binary actions
\begin{equation}
U=
\begin{bmatrix}
u_0^\top & u_1^\top & \cdots & u_{H-1}^\top
\end{bmatrix}^{\top}
\in \{0,1\}^{nH}.
\end{equation}
The number of binary variables is therefore
\begin{equation}
q=nH.
\end{equation}
For example, $n=8$ and $H=8$ gives $q=64$ binary variables, which correspond to a 64-qubit QUBO circuit.

The finite-horizon LQR-style objective is
\begin{equation}
J(U)=
\sum_{k=0}^{H-1}
\left(x_k^\top Qx_k+u_k^\top Ru_k\right)
+x_H^\top Qx_H.
\label{eq:fhcost}
\end{equation}
This objective penalizes unsafe states across the horizon and also penalizes unnecessary defense actions.

\subsection{Exact QUBO Reduction}

Because the system dynamics are linear and the control vector is binary, every future state can be written as an affine function of the stacked control vector:
\begin{equation}
x_k=d_k+M_kU.
\end{equation}
The recursion used in the implementation is initialized by
\begin{equation}
d_0=x_0,\qquad M_0=0,
\end{equation}
and updated as
\begin{equation}
d_{k+1}=Ad_k+Fw_k,
\end{equation}
\begin{equation}
M_{k+1}=AM_k+BE_k,
\end{equation}
where $E_k$ selects the block $u_k$ from the stacked vector $U$.

Substituting $x_k=d_k+M_kU$ into Eq.~\eqref{eq:fhcost} yields a QUBO:
\begin{equation}
\min_{U\in\{0,1\}^{q}}
U^\top H_q U+g_q^\top U.
\label{eq:qubo}
\end{equation}
The matrices are obtained by accumulating state and control contributions:
\begin{equation}
H_q =
\sum_{k=0}^{H} M_k^\top Q M_k
+
\sum_{k=0}^{H-1}E_k^\top R E_k,
\end{equation}
\begin{equation}
g_q =
2\sum_{k=0}^{H}M_k^\top Qd_k.
\end{equation}
This transformation is exact for the finite-horizon linear-quadratic model. No artificial encoding is introduced: the QUBO follows directly from standard control theory.

\section{Optimization Methods}

\subsection{Greedy Baseline}

The Greedy method is included only as a simple baseline. It activates a binary variable when its individual diagonal contribution appears beneficial:
\begin{equation}
U_i=1 \quad \mathrm{if} \quad (H_q)_{ii}+(g_q)_i<0.
\end{equation}
This method ignores coupled interactions and is therefore expected to be weaker than CP-SAT and QAOA.

\subsection{CP-SAT Reference}

CP-SAT is used as the primary classical reference. The QUBO contains quadratic binary terms, so the implementation introduces auxiliary variables $y_{ij}=U_iU_j$ for nonzero pairwise terms. The constraints
\begin{equation}
y_{ij}\leq U_i,\qquad y_{ij}\leq U_j,\qquad y_{ij}\geq U_i+U_j-1
\end{equation}
linearize the product. CP-SAT then minimizes the scaled linearized objective. CP-SAT is not described as an exact proof of global optimality for all large instances because a time limit is used; instead, it is treated as a strong classical reference.

\subsection{QAOA and Ising Mapping}

QAOA requires an Ising Hamiltonian. Using
\begin{equation}
U_i=\frac{1-z_i}{2},\qquad z_i\in\{-1,+1\},
\end{equation}
the QUBO in Eq.~\eqref{eq:qubo} is converted to an Ising form
\begin{equation}
\hat{H}_C=\sum_i h_iZ_i+\sum_{i<j}J_{ij}Z_iZ_j.
\end{equation}
The implementation includes a self-check verifying that the Ising energy plus its constant offset reproduces the original QUBO energy for all bitstrings in a small test instance. This is important because an incorrect QUBO-to-Ising field shift would cause QAOA and CP-SAT to optimize different objectives.

For $p=1$, the QAOA state is
\begin{equation}
|\psi(\gamma,\beta)\rangle =
U_M(\beta)U_C(\gamma)|+\rangle^{\otimes q},
\end{equation}
where
\begin{equation}
U_C(\gamma)=e^{-i\gamma \hat{H}_C}, \qquad
U_M(\beta)=e^{-i\beta\sum_i X_i}.
\end{equation}
The variational angles $(\gamma,\beta)$ are optimized locally using COBYLA. In the validation panel, Local QAOA is run on the full QUBO. The optimized validation angles are then reused for IBM hardware execution.

\subsection{IBM Quantum Execution}

IBM QAOA is executed using Qiskit Runtime SamplerV2 on a real IBM quantum backend. In the reported experiments, the backend selected was \texttt{ibm\_marrakesh}. Each measured bitstring is decoded using Qiskit's little-endian convention and evaluated using the original full QUBO energy. For the larger hardware-scale instances, the QAOA circuit retains the strongest pairwise couplings under a circuit-depth budget; however, all returned bitstrings are scored on the full QUBO. This makes the evaluation conservative: IBM QAOA is graded against the same objective used by CP-SAT, even when the quantum circuit uses a pruned Hamiltonian for hardware feasibility.

\section{Experimental Design}

The experimental design separates algorithmic validation from hardware scalability.

\subsection{Panel 1: Validation}

The validation panel uses $n=4$ nodes and horizon $H=4$, giving
\begin{equation}
q=nH=16
\end{equation}
binary variables. All methods solve the same full QUBO, with no coupling pruning:
\begin{itemize}
    \item Greedy baseline;
    \item CP-SAT reference;
    \item Local QAOA simulation;
    \item IBM QAOA hardware execution.
\end{itemize}
The purpose of this panel is to verify that QAOA solves the correct objective and that IBM hardware produces solutions comparable to the local simulation and CP-SAT reference.

\subsection{Panel 2: Hardware-Scale Demonstration}

The scalability panel evaluates:
\begin{equation}
(n,H)\in\{(8,8),(16,8),(19,8)\},
\end{equation}
corresponding to
\begin{equation}
q\in\{64,128,152\}.
\end{equation}
Only CP-SAT and IBM QAOA are compared, because local full-state QAOA simulation becomes increasingly expensive for dense QUBOs of these sizes. The number of retained quantum couplings is limited to 300, 500, and 700 for 64, 128, and 152 qubits, respectively.

\subsection{Evaluation Metric}

The primary metric is gap-to-CP-SAT:
\begin{equation}
\mathrm{Gap}(\%)=
100
\frac{E_{\mathrm{method}}-E_{\mathrm{CP\text{-}SAT}}}
{|E_{\mathrm{CP\text{-}SAT}}|}.
\label{eq:gap}
\end{equation}
CP-SAT has gap $0\%$ by definition. Raw QUBO energy is reported inside each instance but is not used to compare across different problem sizes.

\section{Results}

\subsection{Validation Panel}

Table~\ref{tab:panel1} presents the validation results on the 16-qubit full QUBO instance. CP-SAT achieves the best objective value and is used as the reference solution. Local QAOA obtains an energy of $-136.7491$, corresponding to a gap of only $0.952\%$ relative to CP-SAT. IBM QAOA achieves an energy of $-136.0254$, corresponding to a $1.476\%$ gap. In contrast, the Greedy baseline performs substantially worse, with a gap exceeding $210\%$.

The close agreement between Local QAOA and IBM QAOA indicates that the variational solution obtained during local optimization transfers successfully to real quantum hardware. Although solution quality degrades as problem size increases, IBM QAOA remains within approximately $8\%$ of the CP-SAT reference even for the 152-qubit instance. The observed increase in gap is consistent with the reduced coupling retention required to satisfy current hardware constraints and circuit-depth limitations.

\begin{table}
\centering
\caption{Panel 1 validation results on the same full 16-qubit QUBO.}
\label{tab:panel1}
\begin{tabular}{lrrrr}
\toprule
Method & QUBO Energy & Gap to CP-SAT (\%) & Runtime (s) & Backend \\
\midrule
Greedy & 152.8038 & 210.676 & 0.000 & Classical \\
CP-SAT & -138.0638 & 0.000 & 0.045 & Classical Reference \\
Local QAOA & -136.7491 & 0.952 & 0.429 & AerSimulator \\
IBM QAOA & -136.0254 & 1.476 & 7.308 & ibm\_marrakesh \\
\bottomrule
\end{tabular}
\end{table}

Figure~\ref{fig:panel1} visualizes the validation results. Both Local QAOA and IBM QAOA remain close to the CP-SAT reference, whereas the Greedy baseline produces significantly poorer solutions.

\begin{figure}
\centering
\includegraphics[width=0.85\textwidth]{panel1_validation.png}
\caption{Validation panel results for the 16-qubit instance. Local QAOA and IBM QAOA achieve solutions within approximately 1--1.5\% of the CP-SAT reference.}
\label{fig:panel1}
\end{figure}

\subsection{Scalability Panel}

Table~\ref{tab:panel2} summarizes the hardware-scale experiments at 64, 128, and 152 qubits. CP-SAT is used as the classical reference, while IBM QAOA is executed on \texttt{ibm\_marrakesh} under a coupling-budget constraint.

At 64 qubits, IBM QAOA achieves a gap of $4.132\%$ relative to CP-SAT. The gap increases to $7.079\%$ at 128 qubits and $7.900\%$ at 152 qubits. Although solution quality degrades as problem size increases and fewer couplings can be retained within the hardware budget, the quantum solutions remain reasonably close to the classical reference.

\begin{table}
\centering
\caption{Panel 2 scalability results. IBM QAOA is evaluated on the full QUBO while the quantum circuit retains only the strongest couplings to satisfy hardware constraints.}
\label{tab:panel2}
\begin{tabular}{rlrrrrl}
\toprule
Qubits & Method & QUBO Energy & Gap (\%) & Retained Couplings (\%) & Runtime (s) & Backend \\
\midrule
64 & CP-SAT & -3603.5012 & 0.000 & 100.00 & 20.045 & Classical \\
64 & IBM QAOA & -3454.6028 & 4.132 & 60.09 & 7.470 & ibm\_marrakesh \\

128 & CP-SAT & -6996.5503 & 0.000 & 100.00 & 20.120 & Classical \\
128 & IBM QAOA & -6501.2573 & 7.079 & 43.45 & 7.827 & ibm\_marrakesh \\

152 & CP-SAT & -8179.5328 & 0.000 & 100.00 & 20.219 & Classical \\
152 & IBM QAOA & -7533.3758 & 7.900 & 46.43 & 32.916 & ibm\_marrakesh \\
\bottomrule
\end{tabular}
\end{table}

The increase in gap from $4.132\%$ at 64 qubits to $7.900\%$ at 152 qubits remains sublinear relative to the increase in problem size. This suggests that the proposed LQR-to-QUBO formulation retains useful optimization structure even when executed through a heavily constrained quantum circuit representation.

The primary observation is that the IBM QAOA gap remains below $8\%$ even at 152 qubits. While these results do not constitute quantum advantage, they demonstrate that the proposed LQR-to-QUBO cyber-defense framework can be executed on contemporary IBM quantum hardware while maintaining solution quality reasonably close to a strong classical optimization baseline.

Figure~\ref{fig:panel2} illustrates the scalability behavior. As the number of qubits increases, the gap to CP-SAT grows gradually rather than catastrophically, suggesting that the methodology scales in a controlled manner under current hardware limitations.

\begin{figure}
\centering
\includegraphics[width=0.85\textwidth]{panel2_scalability.png}
\caption{Scalability results showing the gap to the CP-SAT reference for IBM QAOA at 64, 128, and 152 qubits.}
\label{fig:panel2}
\end{figure}

\section{Discussion}

The validation panel is important because it establishes objective consistency. All methods solve the same 16-qubit full QUBO, and both Local QAOA and IBM QAOA are within $1.5\%$ of CP-SAT. This reduces the risk that the quantum results are an artifact of a mismatched Hamiltonian, incorrect bit ordering, or an unfair evaluation metric.

The scalability panel addresses a different question: whether the workflow remains executable on larger IBM hardware-scale instances. The result is not a proof of quantum advantage. CP-SAT remains a strong classical reference, and QAOA uses a shallow $p=1$ circuit. However, the IBM solution quality remains within approximately $8\%$ of CP-SAT up to 152 qubits. This is meaningful because the quantum hardware is executing a control-derived QUBO rather than an artificial benchmark disconnected from the cyber-defense model.

The use of gap-to-CP-SAT is central to the interpretation. Raw QUBO energies naturally scale with the number of variables and should not be compared directly across problem sizes. By normalizing each method against CP-SAT on the same instance, the evaluation avoids the misleading conclusion that smaller problems are automatically better because they have smaller absolute costs.

The retained-coupling ratio in the scalability panel is also important. At 64 qubits the IBM circuit retains $60.09\%$ of the coupling weight, while at 128 and 152 qubits the retained fractions are $43.45\%$ and $46.43\%$. Thus, the large-scale quantum runs should be interpreted as hardware-feasibility demonstrations under a circuit-depth budget, not as full-depth QAOA optimization of the dense QUBO Hamiltonian.

\section{Limitations}

Several limitations remain. First, CP-SAT is used as a reference under a finite time limit, so it should not be described as a mathematical proof of optimality for the larger instances. Second, QAOA is restricted to $p=1$ to maintain feasible circuit depth on current hardware. Higher-depth QAOA may improve solution quality but would also increase noise sensitivity and transpilation overhead. Third, large-scale IBM runs use coupling truncation for hardware feasibility; although returned bitstrings are evaluated on the full QUBO, the quantum circuit itself optimizes a pruned Hamiltonian. Fourth, the current formulation uses binary defense actions. Continuous defense intensities or multi-level actions would require different encodings. Finally, the present implementation solves a finite-horizon plan; a receding-horizon deployment that repeatedly observes the network, applies the first action, and re-optimizes from the updated state is an important direction for future work.

\section{Conclusion}

This paper presented a hybrid quantum--classical framework for cyber-defense optimization. The method models a network through an extended state vector containing compromise, attacker knowledge, vulnerability exposure, and asset criticality. A finite-horizon LQR-style objective is transformed exactly into a QUBO over binary defense actions. The resulting QUBO is solved using Greedy optimization, CP-SAT, Local QAOA, and IBM QAOA.

On a 16-qubit validation instance, Local QAOA and IBM QAOA achieved gaps of $0.952\%$ and $1.476\%$ relative to CP-SAT, while the Greedy baseline was more than $210\%$ away from the reference. On larger hardware-scale instances, IBM QAOA achieved gaps of $4.132\%$, $7.079\%$, and $7.900\%$ at 64, 128, and 152 qubits, respectively. These results do not establish quantum advantage. Instead, they demonstrate that a control-derived cyber-defense QUBO can be executed on real IBM quantum hardware while preserving solution quality reasonably close to a strong classical benchmark.

Future work will investigate receding-horizon control, improved QAOA parameter transfer, deeper but noise-aware QAOA circuits, error mitigation, and larger cyber-defense topologies with richer operational constraints.

\begin{thebibliography}{13}

\bibitem{astafieva2025optimal}
Astafieva, M., Stepanenko, N.: Optimal cybersecurity management: Global controllability of linear models. In: Workshop on Cybersecurity Providing in Information and Telecommunication Systems (CPITS 2025), CEUR Workshop Proceedings, vol. 3991, pp. 14--25 (2025)

\bibitem{bennett1984quantum}
Bennett, C.H., Brassard, G.: Quantum cryptography: Public key distribution and coin tossing. In: Proc. IEEE International Conference on Computers, Systems and Signal Processing, pp. 175--179 (1984)

\bibitem{garg2022quantum}
Garg, S., Ramakrishna, K.: Quantum machine learning for cyber threat intelligence: A survey. IEEE Access 10, 112305--112318 (2022)

\bibitem{farhi2014qaoa}
Farhi, E., Goldstone, J., Gutmann, S.: A quantum approximate optimization algorithm. arXiv:1411.4028 (2014)

\bibitem{cerezo2021variational}
Cerezo, M., et al.: Variational quantum algorithms. Nature Reviews Physics 3(9), 625--644 (2021)

\bibitem{lubinski2024optimization}
Lubinski, T., Coffrin, C., McGeoch, C., Sathe, P., Apanavicius, J., Bernal Neira, D.E.: Optimization applications as quantum performance benchmarks. arXiv:2302.02278 (2024)

\bibitem{mcclean2018barren}
McClean, J.R., Boixo, S., Smelyanskiy, V.N., Babbush, R., Neven, H.: Barren plateaus in quantum neural network training landscapes. Nature Communications 9, 4812 (2018)

\bibitem{khanafer2016stability}
Khanafer, A., Başar, T., Gharesifard, B.: Stability of epidemic models over directed graphs: A positive systems approach. Automatica 74, 126--134 (2016)

\bibitem{zhu2013game}
Zhu, Q., Başar, T.: Game-theoretic approach to feedback-driven multi-stage moving target defense. In: Decision and Game Theory for Security, pp. 246--263 (2013)

\bibitem{googlecolab}
Google Research: Colaboratory: Google's cloud-based Jupyter environment. \url{https://colab.research.google.com}

\bibitem{ibmquantum}
IBM Quantum: IBM Quantum Platform. \url{https://quantum.ibm.com}

\bibitem{qiskit}
Qiskit contributors: Qiskit: An open-source framework for quantum computing. \url{https://doi.org/10.5281/zenodo.2573505}

\bibitem{ortools}
Perron, L., Furnon, V.: OR-Tools. Google. \url{https://developers.google.com/optimization}

\end{thebibliography}

\end{document}
